% SPDX-License-Identifier: LPPL-1.3c
% Copyright (C) 2026 Christophe Jorssen
% Author and Current Maintainer: Christophe Jorssen <christophe.jorssen@gmail.com>
% LPPL maintenance status: maintained
% This file is part of luafemm. See LICENSE and LICENSES.md for its terms.
\section{Numerical method, verification and limits}
\label{sec:numerics}
The transverse axial-H diffusion formulation and its Bessel benchmarks are
detailed in Section~\ref{sec:currents}; it uses a different scalar unknown
from the vector-potential formulations below.

\subsection{Planar finite elements}
The domain is invariant along $z$. In SI units,
$\mathbf A=A_z\mathbf e_z$ and
$\mathbf B=(\partial_y A_z,-\partial_x A_z)$.
For test functions $v$ vanishing on the Dirichlet boundary, the weak problem is
\[
 \int_\Omega \nu\nabla A_z\cdot\nabla v+\int_{\Gamma_w}(cA_z+g)v
 =\int_\Omega J_zv+
  \int_\Omega H_c\mathbf m\cdot(\partial_yv,-\partial_xv).
\]
Here $\Gamma_w$ contains Neumann and Robin sides, with $c=0$ for Neumann.
The boundary tangent is counterclockwise; $H_t=cA_z+g$. On an edge of
length $\ell$, the boundary mass matrix is
$\ell\bigl(\begin{smallmatrix}2&1\\1&2\end{smallmatrix}\bigr)/6$.
It integrates both $cA_z$ and the affine data $g$ exactly. These terms
are added to the residual and, for $cA_z$, its tangent. Corner constraints
and the all-Neumann compatibility check are described in
Section~\ref{sec:boundaries}.
The potential uses continuous piecewise-affine functions on triangles.
Material laws and sources are constant within each element. With basis
functions $N_i$, the element residual is
\[
 R_i=|T|\left[\nu\nabla A_z\cdot\nabla N_i
 -J_z/3-H_c(m_x\partial_yN_i-m_y\partial_xN_i)\right].
\]
The differential tangent follows from the derivative of the selected
$H(B)$ law: piecewise linear by default, or the FEMM-compatible cubic policy
described in Section~\ref{sec:interchange}. The code scales the entire system by $\mu_0$ for numerical
convenience; this does not change the solution. Dirichlet degrees of
freedom are eliminated. Newton steps use backtracking to decrease the
residual. Linear systems use conjugate gradients and incomplete Cholesky
IC(0); an unsuccessful factorisation is restarted with a positive diagonal
shift of the preconditioner, without changing the system being solved.

The default boundary $A_z=0$ represents a finite truncation of the
exterior. For physical depth $d$, potential levels separated by $\Delta A_z$
enclose flux $d\,\Delta A_z$. Equal level spacing therefore represents equal
flux spacing, not equal geometric separation.

\subsection{Linear algebra, acceptance and recovery}
\label{sec:linear-acceptance}
By default, static Newton corrections and auxiliary force-weight potentials use real
conjugate gradients (CG) with an incomplete Cholesky factor. Harmonic systems
use restarted GMRES with a complex incomplete LU factor, applied through its
real block representation. GMRES does not assume complex symmetry. A known
indefinite constrained operator must not be sent to CG. Gauges and compatible
sources are handled by the physical model before this algebraic stage.

Each solve validates dimensions, sparse indices and finite coefficients and
loads. It equilibrates rows and columns symmetrically and normalises the load;
scaled sums of squares avoid losing a small nonzero source by squaring it.
An exactly zero load returns the zero solution after structural checks and
before factorisation. This shortcut does not certify that an arbitrary matrix
has a unique solution. A singular empty row, a nonfinite quantity or a scaling
operation that loses a nonzero coefficient raises an error.

Let $Kx=b$ denote the assembled linear system after constraints. Before accepting
a solution, the solver recomputes $r=b-Kx$ from the original coefficients and
the returned unknowns in their original units. Neither a shifted factor nor
an iterative recurrence supplies this final residual. With
$c(z)=|\mathop{\rm Re}z|+|\mathop{\rm Im}z|$, the two acceptance measures are
\[
 \rho=\frac{\|r\|_2}{\|b\|_2},\qquad
 \eta=\max_i\frac{c(r_i)}{c(b_i)+\bigl(\max_j c(x_j)\bigr)\sum_j c(K_{ij})}.
\]
The norm in $\rho$ is the Euclidean norm of the real and imaginary components;
for real systems these formulas reduce to their real counterparts. A zero
numerator and denominator contribute zero to $\eta$. Both criteria must pass:
$\rho\leq\tau$ and $\eta\leq\max(\tau,10^{-14})$.
This row-scaled normwise diagnostic limits the effect of unequal row scales
without requiring relative accuracy in an almost-zero solution component. Neither criterion estimates the discretisation error in $\mathbf B$,
material uncertainty or the error of a force integral.

The iteration also checks the true residual at restarts. CG checks recurrence
drift periodically and replaces the recurrence only when needed, preserving
conjugacy otherwise. If conversion back to physical units spoils acceptance,
the solver continues with the recomputed residual within the same budget.
GMRES uses two orthogonalisation passes. The automatic limits are 4000 linear
iterations and a GMRES restart length of 40; exhaustion is an error.

A failed incomplete factor is discarded. At most six attempts use finite,
scaled diagonal shifts, ending with a row-dominance estimate when appropriate;
a diagonal preconditioner is the final fallback. Positive Cholesky pivots and
finite factors are required. \emph{Only the preconditioner is shifted}: the
physical operator, residual and energy remain unchanged. Recovery can make a
factor usable but cannot guarantee convergence within the finite budget.
No numerical factor is reused after a change of operator.

Sorted packed rows are used for iterative products, retaining the original
mesh order. Algebraic ordering and bounded additional fill have been compared
on static, harmonic, periodic, moving-band and biased-response problems.
Their benefit depends on the problem, so automatic solving retains mesh order
and zero-fill factors. Reverse Cuthill--McKee ordering can be requested explicitly with the keys below;
additional-fill experiments remain private.
Validation and the final original-system check add work; a universal speedup
over less checked iteration is not claimed.

\paragraph{Tolerance is not mesh accuracy.}
The existing |/femm/problem/tolerance| controls the outer Newton residual for a
static nonlinear model. The initial residual is bounded below by $10^{-30}$;
the existing absolute stopping floor is $10^{-25}$ in the assembled,
$\mu_0$-scaled system. Automatic static inner corrections use $\tau=10^{-9}$,
independently of that outer target. Harmonic linear solves use the problem's
|tolerance| (strictly between zero and one) as their linear target; the linear kernel has no absolute
small-source cutoff. |max iterations| controls Newton corrections, not the
configurable linear budget. These different meanings matter when interpreting
|\femmstat{residual}|.

This example changes the outer target through an ordinary reusable TikZ style.
The field pic triggers the solve after the geometry has been declared. Its
reported residual is the nonlinear residual; decreasing it does not refine
the displayed coarse mesh. Compare fields on finer meshes separately.

\noindent\begin{minipage}{\linewidth}
\begin{codeexample}[width=0pt]
\tikzset{draft magnetic solve/.style={femm/problem={
  xmin=-65,xmax=65,ymin=-40,ymax=65,
  mesh size=8,tolerance=1e-6}}}
\begin{tikzpicture}[draft magnetic solve,scale=.5]
  \node[femm/u electromagnet={gap=3,gap mesh size=0}] {};
  \pic[femm/lines=9] {femm field};
  \node[below,align=center,font=\small] at (0,-42)
    {Triangles: \femmstat{triangles}\\
     Relative Newton residual: \femmstat{residual}};
\end{tikzpicture}
\end{codeexample}
\end{minipage}

Cached results carry the numerical policy revision. Revision 22 invalidates
older snapshots instead of presenting an unchecked legacy result as newly
verified. A failed solve is not saved as a converged solution. CPU statistics
restored from a cache describe its original calculation, not the loading cost.

\subsection{Choosing and inspecting a linear solve}
\label{sec:linear-controls}
The defaults need no additional keys. The following controls belong to
|/femm/problem| and are read when the picture creates its model, before paths
are captured. They compose with |.style| and |.append style|, obey TeX grouping,
and use the last explicit assignment. An imported FEM problem accepts explicit
numerical overrides. A saved operating point or sliding template inherits its
numerical settings; explicitly supplied numerical keys override that copy without
changing the retained source. Passive ANS answers have no local solver history
and reject numerical overrides. These controls do not change exported geometry,
FEMM node numbers or file units, and are not encoded as FEMM solver settings.

\begin{key}{/femm/problem/linear solver=\meta{choice} (initially auto)}
Choose |auto|, |cg| or |gmres|. Automatic selection uses CG for real positive
static tangents (including the supported lossless fundamental law), and GMRES
for complex harmonic or transverse axial-H diffusion systems. Explicit |gmres|
is also allowed for real static systems. Explicit |cg| is rejected for the
general/complex class, even with zero excitation. Selection follows the operator,
not merely the frequency. Auxiliary force weights always use their own SPD CG
solve; they do not become complex when the magnetic model is harmonic.
\end{key}
\begin{key}{/femm/problem/preconditioner=\meta{choice} (initially auto)}
Choose |auto|, |ic0|, |ilu0| or |diagonal|. Automatic selection starts with IC(0)
for CG and ILU(0) for GMRES. Explicit |ic0| requires CG; explicit |ilu0| requires
GMRES. Choosing a factor does not implicitly change the solver. Both incomplete
factors retain the bounded shift/recovery policy described above. |diagonal|
skips incomplete factorisation. This inexpensive option may need many more
iterations. Force weights inherit |diagonal| when selected; otherwise they use
IC(0) with recovery independently of the magnetic factor choice.
\end{key}
\begin{key}{/femm/problem/linear tolerance=\meta{number or auto} (initially auto)}
A dimensionless numeric literal strictly between zero and one, or |auto|.
The target applies to both original-system acceptance tests above. Automatic
static Newton corrections use $10^{-9}$; harmonic and axial-H systems use the
problem's |tolerance|. An explicit value overrides that inner target without
changing the outer Newton criterion. Circuit-network response columns use the
stricter $\min(0.01\tau,10^{-12})$ to control cancellation in their coupled
solution, whose final residual must meet $\tau$. Auxiliary force weights retain
the measurement's own |tolerance|. Tight algebraic convergence does not certify
field or force accuracy on a coarse mesh.
\end{key}
\begin{key}{/femm/problem/linear max iterations=\meta{integer} (initially 4000)}
A dimensionless integer literal from 1 to 1000000. Each inner solve has this
finite iteration budget, including recurrence replacements and GMRES restarts.
It also applies to force-weight potentials and to each network response column;
the sum reported for a complete magnetic calculation can exceed this number.
Exhaustion raises an error. The separate |max iterations| still bounds Newton
corrections.
\end{key}
\begin{key}{/femm/problem/gmres restart=\meta{integer} (initially 40)}
A dimensionless integer literal from 1 to 4000: the maximum retained GMRES
Krylov basis length before restarting. Increasing it trades memory for possible
convergence improvements. It is validated but unused by CG. It never extends
|linear max iterations|.
\end{key}
\begin{key}{/femm/problem/linear ordering=\meta{choice} (initially mesh)}
Choose |mesh| or |rcm| (reverse Cuthill--McKee). The permutation affects only
algebraic rows and columns, with the returned solution mapped back to its
original mesh. It may change fill behaviour and iteration counts, without
changing the finite-element problem. There is no universal performance gain.
The automatic default remains mesh order and zero-fill incomplete factors.
\end{key}

Here a reusable numerical style is extended and then overridden for two solves
of the same affine-field patch. The exact field is $B_x=0.02$ T,
$B_y=-0.01$ T. The printed iteration counts expose the computational effect;
the two fields should agree. This small case is an illustration, not a speed
benchmark. All six controls are explicit to show their spelling; ordinary
figures can omit them.

\noindent\begin{minipage}{\linewidth}
\begin{codeexample}[width=0pt]
\pgfkeys{/femm/problem/patch solve/.style={
  linear solver=auto,preconditioner=auto,
  linear tolerance=1e-9,linear max iterations=1000},
 /femm/problem/patch solve/.append style={
  gmres restart=20,linear ordering=mesh}}
\foreach \method/\factor/\order in
  {cg/ic0/mesh,gmres/diagonal/rcm}{%
 \begin{tikzpicture}[
  femm/problem={xmin=-10,xmax=10,ymin=-10,ymax=10,
   mesh size=4,patch solve,linear solver=\method,
   preconditioner=\factor,linear ordering=\order},
  femm/boundary={x coefficient=.01,y coefficient=.02}]
  \pic[femm/lines=7] {femm field};
  \draw[gray] (-10,-10) rectangle (10,10);
  \node[below,align=center,font=\small] at (0,-11)
   {\method: \femmstat{linear_iterations} iterations\\
    $|\mathbf B|=\femmB{0}{0}$ T};
 \end{tikzpicture}\quad
}
\end{codeexample}
\end{minipage}

\paragraph{Linear statistics.}
The existing |\femmstat| command accepts the following statistic names for the
current solved magnetic model. Counts and times sum over its
inner solves; maxima are indicated explicitly. Force-weight measurements have
separate solves and do not retroactively change these magnetic statistics.
An unsolved model or an imported ANS answer has no such local history and raises
an error. A native zero-Newton solution reports zero linear work. A successful
zero-load inner solve counts as one solve, with zero iterations and no factor.

\begin{center}\small
\begin{tabular}{lp{.49\linewidth}}
\toprule Name & Meaning \\
\midrule
|assembly_seconds| & Physical assembly CPU seconds (including static backtracking).\\
|linear_solves| & Number of inner solves.\\
|linear_iterations| & Total inner iterations, including GMRES.\\
|linear_cg_solves|, |linear_gmres_solves| & Solves using each iteration method.\\
|linear_diagonal_solves| & Solves using a diagonal factor, explicit or fallback.\\
|linear_restarts| & Iteration restart count.\\
|linear_retries| & Incomplete-factor retry count.\\
|linear_nonzeros| & Maximum stored factor nonzeros.\\
|linear_shift| & Maximum dimensionless equilibrated factor shift.\\
|linear_residual| & Maximum true relative residual $\rho$.\\
|linear_backward_error| & Maximum row-scaled backward error $\eta$.\\
|linear_finalize_seconds| & Sparse preparation CPU seconds.\\
|linear_factor_seconds| & Factor construction CPU seconds.\\
|linear_iteration_seconds| & Iteration and acceptance CPU seconds.\\
\bottomrule
\end{tabular}
\end{center}
These values describe linear subproblems; for voltage-driven networks the final
coupled acceptance remains a separate check. |residual| keeps its model-specific
meaning. Numerical options participate in the native solution-cache identity:
changing any of them invalidates the solution, while permitting compatible mesh
reuse. Timings restored from cache describe the original calculation, not cache
loading. Factors themselves are never saved or reused across changed operators.

\subsection{Constrained meshing and predicates}
The mesher splits crossings, T junctions and collinear overlaps of the
polygonal constraints. It samples their edges and seeds the interior
with triangular lattices, using the finest local size where requested.
Incremental Delaunay insertion is followed by constrained-segment cavity
retriangulation and Lawson edge flips.

Topological decisions use orientation and incircle predicates. A
floating-point filter certifies well-separated signs; uncertain cases
are recomputed using error-free sums/products and floating-point
expansions in Lua. The result has the exact sign of the determinant of
the supplied binary coordinates, including exact zero, provided arithmetic
neither overflows nor underflows. Normalising coordinates before meshing
keeps the intended working range moderate. This does not make the input
geometry exact before its TeX rounding, nor make intersection construction
exact: intersection coordinates still use floating-point arithmetic.

The native staging-mesh audit checks positive areas, rectangular coverage, preservation
of every constraint and vertex, at most two triangles per edge, exposed
edges only on the tagged exterior, Euler's relation $V-E+F=1$, and local
Delaunay legality of every unconstrained flippable edge. Failure stops
the calculation. The vertex limit is 40000 vertices for Delaunay and
150000 grid vertices.

Triangle was studied as an algorithmic reference; its source and binary
are not embedded. This is not a full Triangle port. Optional segment-first
Delaunay refinement enforces angle and area targets, or raises an error
when its budget or geometric precision is exhausted (Section~\ref{sec:mesh-quality}).
There is no specialised acute-corner protection or guaranteed smooth
transition between mesh sizes. With both quality targets disabled, narrow
features and acute corners can create badly conditioned triangles.

Conservative spatial queries accelerate segment intersections, vertex-on-segment
searches, lattice clearance and refinement encroachment tests. They retain the
final predicates and ascending candidate order; they do not replace geometric
tests with grid membership. Exhaustive-search comparisons on translated, scaled
and reversed contours verify identical meshes on the regression corpus.
Manufactured fields check shared periodic corners and moving-ring references
after tolerance-based vertex merging. These checks do not prove correctness for
every geometry or implement specialised corner protection.
Monitor |min_angle| and perform field-convergence studies; triangle quality
alone is not a physical error estimate.

\subsection{What is tested}
The release checks analytic affine fields for Dirichlet, Neumann, Robin
and mixed sides on both meshers (including nonlinear permanent magnets),
manufactured Poisson and Robin convergence, all-Neumann compatibility,
nonlinear saturation, current and magnetisation reversal, an exact
magnet/air interface patch, and a magnetised disk against the infinite-domain
solution. Geometry tests cover intersections, shared edges, small gaps,
concavity, scaling, translation and compound-region fill rules. Circular-shell
shielding is checked against an analytic cavity field with exact exterior
Dirichlet data. An independent Python rational oracle
checks 4530 orientation, incircle and diametral-disk signs; 49 stress
configurations exercise the base mesher. Quality tests independently check
angles and areas, scaled and translated inputs, compound cavities, refined
boundary tags, affine Robin fields, cache reuse and explicit budget failures.
Python is a development-test dependency only.

All 282 FEMM 5 catalogue records are validated. Profiles are tested for Cartesian
and local components, arc length, reversed traversal, normal offsets,
out-of-domain policy and retained model lifetime. The same nonlinear
scene is compiled under plain LuaTeX, LuaLaTeX and ConTeXt MkIV; computed
sample values and node counts must agree within $10^{-10}$. Load guards,
catcode restoration, transformed paths and PGFPlots legends are exercised.

For the original mixed alnico/illustrative-iron U with 1 mm gaps, sample
results document sensitivity rather than absolute accuracy:
\begin{center}\small
\begin{tabular}{rrrrr}
\toprule Global/gap step (mm) & Domain (mm$^2$) & Nodes & Left $B_y$ (T) & Right $B_y$ (T)\\
\midrule
2 / .5 & $200\times185$ & 11376 & .782876 & -.782907\\
1.5 / .25 & $200\times185$ & 20409 & .777974 & -.778597\\
2 / .5 & $280\times265$ & 22293 & .775140 & -.775140\\
\bottomrule
\end{tabular}\end{center}
The refinements change the gap field by about 0.6 percent, and the larger
domain by about one percent. Material-model error is a separate issue.
The interchange suite additionally loads Lua-written ANS files through xfemm's
independent disk postprocessor and compares $A$, unsmoothed $B$ and $H$.
Independent solves of the exported U-coil and alnico models agree to within
2 percent in the tested gap-field samples at the two draft mesh resolutions.
This is a bounded validation result, not a universal accuracy guarantee;
see Section~\ref{sec:interchange} and the reproducible reports in |build/interop|.

Imported domains are carved from a constrained staging mesh. Their audit also
checks active area, protected exposed edges and the Euler characteristic of the
active cell complex; shared child boundaries are counted as a union. Exclusions
are never silently replaced by air. Unanchored disconnected components have
independent gauges and current-balance checks.

\subsection{Supported physics and practical diagnostics}
The static solver is planar or axisymmetric, with a linear symmetric permeability
tensor or a monotone tabulated law and prescribed permanent-magnet direction.
Section~\ref{sec:axisymmetric} gives the separate cylindrical formulation.
The separate linear harmonic solver in Section~\ref{sec:harmonic} adds
single-frequency diffusion and eddy currents. Section~\ref{sec:biased-harmonic}
extends it to a fixed differential tensor around a nonlinear DC bias.
Real backward-Euler time diffusion and controlled motion are described in
Section~\ref{sec:time}. Hysteresis history and irreversible demagnetisation
remain unsupported. Circular P1 moving bands
are described in Section~\ref{sec:sliding}; the independent full-circle Fourier
operator is described in Section~\ref{sec:machine-studies}. Forces and circular
open-boundary approximations are available as described in their chapters.

If Newton fails, inspect the material data, extreme element aspect ratios
and current scale before raising |max iterations|. For a lost or collapsed
geometric feature, compare its size with PGF capture precision and the
merging tolerance. For a graph with jumps, inspect the elementwise field
and sampling side before requesting graphical smoothing. The log identifies
the failing stage; a solver failure is not turned into a plausible plot.

\subsection{Primary references}
The development was based on studying the source code of FEMM/xfemm
and Triangle, then reimplementing selected formulations and algorithms
in Lua/LuaTeX. No native solver or mesher is embedded. The following primary references
also describe the surrounding algorithms and interfaces:
\begin{itemize}
\item \href{https://www.femm.info/wiki/Documentation}{FEMM documentation},
  planar magnetostatics and material conventions; FEMM revision
  \nolinkurl{7d9e8ed90772ddb939a3256f976e231bcb5cca97}.
\item xfemm revision \nolinkurl{34ac6cc586a346066984526620dfb67383c8e3c5},
  \nolinkurl{cfemm/fsolver/static2d.cpp} and
  \nolinkurl{cfemm/libfemm/CMaterialProp.cpp}.
\item Jonathan Shewchuk,
  \href{https://www.cs.cmu.edu/~quake/robust.html}{Robust Adaptive Floating-Point Geometric Predicates}
  and \href{https://www.cs.cmu.edu/~quake/triangle.research.html}{Triangle algorithm references}.
\item The \href{https://tikz.dev}{PGF/TikZ manual} and the installed
  PGFPlots manual for path transformations, key handlers and axes.
\end{itemize}

\subsection{Circular moving-band operator}
Section~\ref{sec:sliding} describes retained bulk meshes, annular reconnection
and prescribed-weight Maxwell stress. The operator remains planar P1 at each
static angle; exact orientations and mesh coverage are checked before solving.
It does not reuse the previous position's nonlinear field as a frozen material
law. Coenergy differences and mesh/angle sensitivity remain necessary checks.
